% Section 2.2, from lectures/L02/appendix/b_contact_networks.md.
%
% Two corrections, made in the source too: B.3 now shows the truth table its text calls "OR-grindens
% sanningstabell" (the lecture had the sentence but no table), and B.6 speaks of four rules, as its
% table has four rows, where the lecture said three.

\section{Kontaktnät}
\label{c2:app:b}

\subsection{Kontakter och symboler}
\label{c2:b1}
Långt innan det fanns IC-kretsar byggdes styrlogik av kontakter: tryckknappar, gränslägesbrytare
och reläer, kopplade så att en lampa, en motor eller en ventil fick ström precis när den skulle. Så
fungerar fortfarande en stor del av säkerhetskretsarna i industrin, och det är samma logik som en
PLC programmeras med i dag (\secref{c2:b8}).

En \textbf{kontakt} har två lägen, sluten eller öppen, och är därför en logisk variabel. Det finns
två sorter:
\begin{itemize}
  \item En \textbf{slutande kontakt} (NO, \emph{normally open}) är öppen i vila och sluts när den
    påverkas. Tryck på knappen, så leder den.
  \item En \textbf{brytande kontakt} (NC, \emph{normally closed}) är sluten i vila och bryts när
    den påverkas. Tryck på knappen, så slutar den leda.
\end{itemize}

En tryckknapp innehåller ofta flera \textbf{kontaktelement} som påverkas samtidigt, till exempel
ett slutande och ett brytande. Kontaktelementens anslutningar har nummer enligt en standard: den
sista siffran är \textbf{3--4 för ett slutande} element och \textbf{1--2 för ett brytande}, och den
första siffran är elementets platsnummer. Ett slutande element kan alltså heta 13--14 och ett
brytande 21--22.

Kontakterna ritas i ett \textbf{strömvägsschema}: vertikala strömvägar mellan en matningsskena
upptill (+24~V) och en återledare nedtill (0~V), med kontakterna ritade vertikalt i vägen och det
som ska styras, en lampa eller en reläspole, längst ned.

\bookfigure{l02_contact_symbols}{Symbolerna för tryckknappar, reläkontakter, signallampa och
reläspole, med anslutningsnummer.}{c2:fig:contact_symbols}

\noindent Symbolerna ritas alltid i \textbf{viloläge}: den slutande kontakten öppen, den brytande
sluten, som när ingen rör något. Den streckade linjen till vänster är tryckknappens påverkansdon;
en reläkontakt har ingen sådan, eftersom den påverkas av sin spole. Varje komponent har en
beteckning: S för tryckknappar och andra manöverdon, H för lampor och K för reläer.

\subsubsection{Kontakten som logisk variabel}
Översättningen till logik är:
\begin{itemize}
  \item en \textbf{nedtryckt knapp} är \code{1}, en släppt knapp \code{0};
  \item en \textbf{tänd lampa} är \code{1}, en släckt lampa \code{0};
  \item \textbf{lampan lyser när det finns en sluten väg från +24~V till 0~V.}
\end{itemize}
Den sista punkten är hela bron mellan kontakter och logik. Allt i resten av det här avsnittet
följer av den.

\subsection{Seriekoppling = AND}
\label{c2:b2}
Två slutande kontakter efter varandra i samma strömväg. Strömmen måste passera \textbf{båda}, så
lampan lyser bara om S1 \textbf{och} S2 är nedtryckta:

\begin{textdrawing}
H1 = S1 · S2
\end{textdrawing}

\begin{narrowtable}{@{}ccc@{}}
  \thead{S1} & \thead{S2} & \thead{H1} \\
  \midrule
  0 & 0 & 0 \\
  0 & 1 & 0 \\
  1 & 0 & 0 \\
  1 & 1 & 1 \\
\end{narrowtable}

\noindent Det är AND-grindens sanningstabell. Ett vardagsexempel är \textbf{tvåhandsmanövern} på
en press: den går bara om båda knapparna trycks samtidigt, så att båda händerna garanterat är
utanför.

\subsection{Parallellkoppling = OR}
\label{c2:b3}
Två slutande kontakter bredvid varandra, var och en en egen väg mellan samma två punkter. Det
räcker att \textbf{en} av vägarna är sluten, så lampan lyser om S1 \textbf{eller} S2 är nedtryckt,
eller båda:

\begin{console}
H1 = S1 + S2
\end{console}

\begin{narrowtable}{@{}ccc@{}}
  \thead{S1} & \thead{S2} & \thead{H1} \\
  \midrule
  0 & 0 & 0 \\
  0 & 1 & 1 \\
  1 & 0 & 1 \\
  1 & 1 & 1 \\
\end{narrowtable}

\noindent Det är OR-grindens sanningstabell. Ett vardagsexempel är en ringklocka med en knapp vid
framdörren och en vid bakdörren: vilken som helst av dem får den att ringa.

\subsection{Brytande kontakt = NOT}
\label{c2:b4}
En brytande kontakt ensam i strömvägen. Lampan lyser \textbf{när knappen inte är nedtryckt}, och
slocknar när den trycks:

\begin{console}
H1 = S1'
\end{console}

\noindent Det är NOT. Ett vardagsexempel är en stoppknapp: den ska bryta, inte sluta, en krets, och
varför den ska vara brytande just av säkerhetsskäl återkommer i \secref{c2:b7}.

Här är de tre grundkopplingarna i samma figur:

\bookfigure{l02_contact_basics}{AND, OR och NOT som strömvägar mellan +24~V och
0~V.}{c2:fig:contact_basics}

\subsection{Sammansatta kontaktnät}
\label{c2:b5}
Med serie, parallell och brytande kontakter går det att bygga varje logisk funktion.

\subsubsection{NAND och NOR}
De Morgans lagar från \secref{c2:a5} visar hur:
\begin{itemize}
  \item NAND: \mbox{\code{(S1 · S2)' = S1' + S2'}}. Två \textbf{brytande} kontakter
    \textbf{parallellt}: lampan slocknar bara om båda vägarna bryts, alltså om båda knapparna
    trycks.
  \item NOR: \mbox{\code{(S1 + S2)' = S1' · S2'}}. Två \textbf{brytande} kontakter \textbf{i serie}:
    lampan slocknar så snart någon av knapparna trycks.
\end{itemize}

\bookfigure{l02_contact_nand_nor}{NAND som två brytande kontakter parallellt, och NOR som två
brytande kontakter i serie.}{c2:fig:contact_nand_nor}

\noindent Det är De Morgan uttryckt i koppar: inverterade ingångar (brytande kontakter) och bytt
operation (parallell i stället för serie).

\subsubsection{XOR och XNOR: trappkopplingen}
En lampa i en trappa ska kunna tändas och släckas från både övervåningen och nedervåningen. Varje
tryck på vilken knapp som helst ska ändra lampans läge. Med två knappar är det XOR:

\begin{textdrawing}
H1 = S1 · S2' + S1' · S2
\end{textdrawing}

\noindent Varje term är en strömväg med två kontakter i serie, och de två vägarna ligger
parallellt (figur~\ref{c2:fig:contact_xor}). Termen \mbox{\code{S1 · S2'}} behöver S1:s
\textbf{slutande} kontakt och S2:s \textbf{brytande}; termen \mbox{\code{S1' · S2}} tvärtom.
\textbf{Varje knapp används alltså med båda sina kontaktelement}, och det är skälet till att
labbstationens knappar har både slutande och brytande kontakter.

\bookfigure[tbp]{l02_contact_xor}{XOR och XNOR som kontaktnät.}{c2:fig:contact_xor}

Kontakterna med samma beteckning, till exempel de två som heter S1, sitter i samma knapp
och påverkas samtidigt. I schemat ritas de där de gör nytta i strömvägen, inte bredvid varandra.

(En riktig trappbelysning använder vippströmbrytare, som stannar i det läge man lämnar dem, i
stället för tryckknappar. Logiken är densamma.)

\subsection{Från uttryck till kontaktnät och tillbaka}
\label{c2:b6}
Översättningen går åt båda hållen med fyra regler:

\begin{narrowtable}{@{}ll@{}}
  \thead{Uttryck} & \thead{Kontaktnät} \\
  \midrule
  en variabel \code{S1} & en slutande kontakt S1 \\
  en inverterad variabel \code{S1'} & en brytande kontakt S1 \\
  en produkt, \code{AB} & kontakterna i \textbf{serie} \\
  en summa, \code{A + B} & kontakterna \textbf{parallellt} \\
\end{narrowtable}

\noindent Parenteser blir grupper: \mbox{\code{(S2 + S3) · S4'}} är en parallellkoppling av S2 och
S3, i serie med en brytande kontakt S4.

\textbf{Från uttryck till nät:} börja med den yttersta operationen och arbeta dig inåt. Är
uttrycket en summa av termer, rita en parallell gren per term; är det en produkt, rita faktorerna
i serie.

\textbf{Från nät till uttryck:} börja med de innersta grupperna, de minsta serie- och
parallellkopplingarna, och arbeta dig utåt. Kapitel~\ref{c3:ch} går igenom det i detalj.

\subsubsection{Bromsassistenten som kontaktnät}
Uttrycket från \secref{c2:a10}, \mbox{\code{B = D + (S + R)F'}}, blir med tryckknappar i stället för
signaler (S1 för D, S2 för S, S3 för R och S4 för F) och en lampa H1 som bromsljus:

\begin{textdrawing}
H1 = S1 + (S2 + S3) · S4'
\end{textdrawing}

\noindent Den yttersta operationen är en summa: två parallella grenar. Den vänstra är S1 ensam. Den
högra är en produkt: parallellkopplingen av S2 och S3, i serie med S4:s brytande kontakt.

\bookfigure[tbp]{l02_contact_adas}{Bromsassistenten som kontaktnät.}{c2:fig:contact_adas}

Säkerhetsargumentet syns direkt i ritningen (figur~\ref{c2:fig:contact_adas}): förarens knapp S1
har en egen väg till lampan, och felsignalen S4 kan bara bryta assistanssystemets gren. Ingen
kombination av de andra knapparna kan hindra S1 från att tända H1.

\subsection{Reläet}
\label{c2:b7}
Ett \textbf{relä} är en elektriskt styrd strömbrytare: en \textbf{spole} som, när den får ström,
drar till sig ett ankare som påverkar reläets \textbf{kontakter}. En reläkontakt kan vara slutande
eller brytande, precis som en tryckknapps. Spolen och kontakterna har samma beteckning, till
exempel K1, men ritas där de hör hemma i schemat: spolen längst ned i en strömväg, kontakterna i
andra strömvägar.

Två egenskaper gör reläet användbart:
\begin{itemize}
  \item \textbf{En liten ström styr en stor.} Spolen kan drivas av en svag signal, till exempel från
    en mikrodator, medan kontakterna kan koppla en motor eller en lampa på 24~V eller 230~V. Det
    återkommer i kapitel~\ref{c11:ch}, när mikrodatorn ska styra något i verkligheten.
  \item \textbf{En signal kan användas många gånger.} Ett relä kan ha flera kontakter, så en logisk
    variabel, \enquote{K1 har dragit}, kan användas i så många strömvägar som behövs, slutande i
    några och brytande i andra.
\end{itemize}

\subsubsection{Självhållning: en krets som minns}
Den viktigaste reläkopplingen av alla är \textbf{självhållningen}, som finns i nästan varje maskin
med en start- och en stoppknapp:

\bookfigure{l02_relay_self_hold}{Självhållning med startknappen S1, stoppknappen S2 och reläet
K1.}{c2:fig:relay_self_hold}

\noindent Följ vad som händer:
\begin{enumerate}
  \item I vila är S1 och reläkontakten K1 öppna. Spolen får ingen ström, och lampan är släckt.
  \item \textbf{Tryck på start, S1.} Strömväg 1 sluts genom S2 (brytande, alltså sluten) och S1.
    Spolen K1 får ström och drar. Då sluts båda K1-kontakterna: lampan tänds, och K1-kontakten
    parallellt med S1 sluts.
  \item \textbf{Släpp S1.} Strömmen går nu genom K1:s egen kontakt i stället, så spolen håller sig
    själv strömsatt. Lampan fortsätter lysa.
  \item \textbf{Tryck på stopp, S2.} Den brytande kontakten bryter strömväg 1, spolen släpper, och
    båda K1-kontakterna öppnar. När S2 släpps är S1 och K1 öppna igen, och kretsen är tillbaka i
    vila.
\end{enumerate}

Lägg märke till vad som hände i steg 3: \textbf{ingångarna är desamma som i steg 1}, båda knapparna
släppta, men lampan lyser. Utgången beror inte bara på ingångarna nu, utan också på vad som har
hänt tidigare. Kretsen \textbf{minns} att den startades.

Det är inte ett kombinatoriskt nät längre. Ett nät där en utgång återkopplas till sin egen ingång,
och som därför har minne, kallas ett \textbf{sekvensnät}, och det är grunden för allt minne i en
dator, från ett enda register till hela arbetsminnet. Kapitel~\ref{c6:ch} bygger samma sak med
grindar i stället för reläer.

Stoppknappen är brytande av ett säkerhetsskäl: går en ledning till den av, bryts kretsen och
maskinen stannar. Med en slutande stoppknapp skulle en avbruten ledning i stället göra det omöjligt
att stoppa maskinen.

\subsection{Kontaktnät och PLC:ns stegdiagram}
\label{c2:b8}
I dag byggs styrlogik sällan av reläer, utan programmeras i en \textbf{PLC}, en industridator gjord
för att styra maskiner. Det vanligaste språket för att programmera en PLC, \textbf{stegdiagram}
(LD, \emph{ladder diagram}), är ritat för att se ut som ett kontaktnät vänt på sidan:
\begin{itemize}
  \item strömvägarna ligger ned, som stegpinnar mellan två lodräta skenor;
  \item en slutande kontakt ritas \mbox{\code{-| |-}} och en brytande \code{-|/|-};
  \item utgången, en spole eller en lampa, ritas \mbox{\code{-( )-}} längst till höger.
\end{itemize}

\bookfigure{l02_ladder}{Självhållningen som stegdiagram.}{c2:fig:ladder}

\noindent Allt i det här avsnittet gäller oförändrat: serie är AND, parallell är OR, en brytande
kontakt är NOT, och självhållningen fungerar på samma sätt. Den som kan läsa ett kontaktnät kan
läsa ett stegdiagram.

\subsection{Förbered Labb 1}
\label{c2:b9}
I Labb~1 (bilaga~\ref{app:lab1}) bygger du kontaktnät med två tryckknappar, en lampa och ett
aggregat på 24~V. Innan labben ska du ha gjort förberedelseuppgifterna F1--F4 i
labbhandledningen. De går ut på det som det här avsnittet har visat:
\begin{itemize}
  \item skriva sanningstabeller för AND, OR, NOT, NAND, NOR, XOR och XNOR med två knappar;
  \item rita strömvägsscheman för dem, med rätt kontaktsort för varje variabel;
  \item ta fram ett uttryck ur en sanningstabell, och förenkla det;
  \item analysera ett färdigt kontaktnät, som i \secref{c2:b6}.
\end{itemize}

\textbf{Rita alltid schemat innan du kopplar.} Ett schema på papper går att kontrollera mot
sanningstabellen på en minut. En felkopplad sladd på stationen tar mycket längre tid att hitta.
